\subsection{同态复形}

设 (C, d) 和 (C', d') 是两个 A-复形。考虑分次 k-模 Homgr\_A(C, C')（II，第 174–175 页）：对于 \(n \in \mathbf{Z}\) ，Homgr\_A(C, C')\_n 是从 C 到 C' 的、次数为 \(n\) 的 A-线性分次映射所成的 k-模；换言之，Homgr\_A(C, C') 与 A-模典范地等同

\[
\bigoplus_ {n \in \mathbf {Z}} \prod_ {p \in \mathbf {Z}} \operatorname{Hom} _ {\mathrm{A}} \left(\mathrm{C} _ {p}, \mathrm{C} _ {p + n} ^ {\prime}\right) = \bigoplus_ {n \in \mathbf {Z}} \prod_ {p + q = n} \operatorname{Hom} _ {\mathrm{A}} \left(\mathrm{C} _ {p}, \mathrm{C} ^ {\prime q}\right).
\]

我们来定义 k-线性映射：

\[
\mathrm{D} _ {n}: \operatorname{Homgr} _ {\mathbf {A}} (\mathrm{C}, \mathrm{C} ^ {\prime}) _ {n} \rightarrow \operatorname{Homgr} _ {\mathbf {A}} (\mathrm{C}, \mathrm{C} ^ {\prime}) _ {n - 1}, \quad n \in \mathbf {Z},
\]

如下

\[
\mathbf {D} _ {n} (f) = d ^ {\prime} \circ f - (- 1) ^ {n} f \circ d;\tag{1}
\]

我们有

\[
\begin{array}{r l} \mathrm{D} _ {n - 1} \circ \mathrm{D} _ {n} (f) = & \mathrm{D} _ {n - 1} \left(d ^ {\prime} \circ f - (- 1) ^ {n} f \circ d\right) = d ^ {\prime} \circ d ^ {\prime} \circ f - (- 1) ^ {n} d ^ {\prime} \circ f \circ d \\ & - (- 1) ^ {n - 1} d ^ {\prime} \circ f \circ d - f \circ d \circ d = 0. \end{array}
\]

则 (Homgr \(_{A}\) (C, C'), D) 是一个 \(k\) -模复形，称为从 C 到 C' 的同态复形。

例如， \(\operatorname{Homgr}_{\mathbf{A}}(\mathbf{A}, \mathbf{C}')\) 与 \(\mathbf{C}'\) 典范地等同。我们还注意到，对于每一个 \(n \in \mathbf{Z}\) ，有 \(\operatorname{Homgr}_{\mathbf{A}}(\mathbf{C}, \mathbf{C}')\) ( \(n\) ) = \(\operatorname{Homgr}_{\mathbf{A}}(\mathbf{C}, \mathbf{C}'(n))\) .

\(Z_n(\text{Homgr}_A(C, C'))\) 的元素是这样的分次同态 \(f\) ：次数为（下降） \(n\) ，从 \(C\) 到 \(C'\) ，且满足 \(d' \circ f = (-1)^n f \circ d\) ，即从 \(C\) 到 \(C'(n)\) 的复形态射，或再次为从 \(C(p)\) 到 \(C'(p + n)\) （对任意固定的 \(p\) ）的复形态射。它们称为（下降）次数为 \(n\) 的、从 \(C\) 到 \(C'\) 的复形态射；若 \(f, g \in Z_n(\text{Homgr}_A(C, C'))\) 和 \(s \in \text{Homgr}_A(C, C')_{n+1}\) ，则条件 \(g - f = Ds\) 意味着 \(s\) 是连接态射 \(f\) 和 \(g\) （从 \(C\) 到 \(C'(n)\) ）的同伦，从而 \(H_n(\text{Homgr}_A(C, C'))\) 是 \(k\) -模，它由（下降）次数为 \(n\) 的、从 \(C\) 到 \(C'\) 的态射的同伦类组成.

设 \(\alpha\in\mathrm{H}_{n}(\mathrm{Homgr}_{\mathrm{A}}(\mathrm{C},\mathrm{C}^{\prime}))\) 和 \(p\in\mathbf{Z}\) 。令 \(\alpha\) 由 \(f\in\mathbf{Z}_{n}(\mathrm{Homgr}_{\mathrm{A}}(\mathrm{C},\mathrm{C}^{\prime}))\) 代表；则 \(f\) 是从 \(\mathbf{C}\) 到 \(\mathbf{C}^{\prime}(n)\) 的复形态射，因此 \(\mathbf{H}_{p}(f)\) 是从 \(\mathbf{H}_{p}(\mathbf{C})\) 到 \(\mathbf{H}_{p}(\mathbf{C}^{\prime}(n))=\mathbf{H}_{p+n}(\mathbf{C}^{\prime})\) 的同态；由于 \(\mathbf{H}_{p}(f)\) 仅依赖于同伦类 \(\alpha\) （ \(f(\mathbf{X},\mathbf{p.~33,~prop.~3})\) ） ，我们推出一个典范的 \(k\) -模同态

\[
\mathrm{H} _ {n} \left(\operatorname{Homgr} _ {\mathrm{A}} \left(\mathrm{C}, \mathrm{C} ^ {\prime}\right)\right)\rightarrow \operatorname{Hom} _ {\mathrm{A}} \left(\mathrm{H} _ {p} (\mathrm{C}), \mathrm{H} _ {p + n} \left(\mathrm{C} ^ {\prime}\right)\right),
\]

由此得到一个次数为0的分次k-线性映射，称为典范映射

\[
\lambda (\mathrm{C}, \mathrm{C} ^ {\prime}): \mathrm{H} (\operatorname{Homgr} _ {\mathbf {A}} (\mathrm{C}, \mathrm{C} ^ {\prime})) \rightarrow \dot {\operatorname{Homgr}} _ {\mathbf {A}} (\mathrm{H} (\mathrm{C}), \mathrm{H} (\mathrm{C} ^ {\prime})).
\]

\(\lambda(C, C')\) 的分次分量通常记作：

\[
\lambda^ {n} (\mathrm{C}, \mathrm{C} ^ {\prime}): \mathrm{H} ^ {n} \left(\operatorname{Homgr} _ {\mathrm{A}} \left(\mathrm{C}, \mathrm{C} ^ {\prime}\right)\right)\rightarrow \prod_ {p + q = n} \operatorname{Hom} _ {\mathrm{A}} \left(\mathrm{H} _ {p} (\mathrm{C}), \mathrm{H} ^ {q} \left(\mathrm{C} ^ {\prime}\right)\right).
\]

命题1. — 若 C 右零且 C' 左零，则 Homgr \(_{A}\) (C, C') 左零，且典范 k-线性映射

\[
\lambda^ {0} (\mathrm{C}, \mathrm{C} ^ {\prime}): \mathrm{H} ^ {0} \left(\operatorname{Homgr} _ {\mathrm{A}} (\mathrm{C}, \mathrm{C} ^ {\prime})\right)\rightarrow \operatorname{Hom} _ {\mathrm{A}} \left(\mathrm{H} _ {0} (\mathrm{C}), \mathrm{H} ^ {0} \left(\mathrm{C} ^ {\prime}\right)\right)
\]

是双射。

我们有正合序列

\[
0 \longrightarrow \mathrm{H} ^ {0} (\mathrm{C} ^ {\prime}) \stackrel {i} {\longrightarrow} \mathrm{C} ^ {\prime 0} \stackrel {d ^ {\prime 0}} {\longrightarrow} \mathrm{C} ^ {\prime 1}
\]

\[
\mathrm{C} _ {1} \xrightarrow {d _ {1}} \mathrm{C} _ {0} \xrightarrow {p} \mathrm{H} _ {0} (\mathrm{C}) \longrightarrow 0.
\]

另一方面， \(\mathrm{Homgr}_{\mathbf{A}}^{0}(\mathbf{C},\mathbf{C}^{\prime})\) 等同于 \(\mathrm{Hom}_{\mathbf{A}}(\mathbf{C}_{0},\mathbf{C}^{\prime 0})\) ，而 \(\mathbf{Z}^{0}(\mathrm{Homgr}_{\mathbf{A}}(\mathbf{C},\mathbf{C}^{\prime}))\) 则等同于这样的 \(f:\mathbf{C}_{0}\to \mathbf{C}^{\prime 0}\) 所成的集合：它使得 \(d^{\prime 0}\circ f = 0\) , \(f\circ d_{1} = 0\) ； \(\mathbf{B}^{0}(\mathrm{Homgr}_{\mathbf{A}}(\mathbf{C},\mathbf{C}^{\prime}))\) 为零；最后，映射 \(\lambda^0\) 把 \(f\) 模 \(\{0\}\) 的类对应到同态 \(\varphi :\mathrm{H}_0(\mathrm{C})\to \mathrm{H}^0 (\mathrm{C}')\) （它使得 \(f = i\circ \varphi \circ p\) ），由此得出该命题。

设 \(u: \tilde{\mathbf{C}} \to \mathbf{C}\) 和 \(u': \mathbf{C}' \to \tilde{\mathbf{C}}'\) 是复形的态射；则典范同态 \(\operatorname{Homgr}_{\mathbf{A}}(u, u') : \operatorname{Homgr}_{\mathbf{A}}(\mathbf{C}, \mathbf{C}') \to \operatorname{Homgr}_{\mathbf{A}}(\tilde{\mathbf{C}}, \tilde{\mathbf{C}}')\) ，定义为 \(f\mapsto u^{\prime}\circ f\circ u\) 是复形的态射，这由公式（1）立即得出。此外，下面的图表是交换的

\[
\begin{array}{c} \mathrm{H} (\mathrm{Homgr} _ {\mathsf {A}} (\mathrm{C}, \mathrm{C} ^ {\prime})) \xrightarrow {\lambda (\mathrm{C} , \mathrm{C} ^ {\prime})} \mathrm{Homgr} _ {\mathsf {A}} (\mathrm{H} (\mathrm{C}), \mathrm{H} (\mathrm{C} ^ {\prime})) \\ \mathrm{H} (\mathrm{Homgr} _ {\mathsf {A}} (u, u ^ {\prime})) \Big \downarrow \\ \mathrm{H} (\mathrm{Homgr} _ {\mathsf {A}} (\tilde {\mathsf {C}}, \tilde {\mathsf {C}} ^ {\prime})) \xrightarrow {\lambda (\tilde {\mathsf {C}} , \tilde {\mathsf {C}} ^ {\prime})} \mathrm{Homgr} _ {\mathsf {A}} (\mathrm{H} (\tilde {\mathsf {C}}), \mathrm{H} (\tilde {\mathsf {C}} ^ {\prime}))  . \end{array}
\]

命题2. --- a) 设 \(C' \xrightarrow{u} C \xrightarrow{v} C''\) 是 A-复形的正合序列， \(P\) 是投射复形，E 是内射复形（X，第 25 页）。则序列

\[
\operatorname{Homgr} _ {\mathbf {A}} (\mathrm{P}, \mathrm{C} ^ {\prime}) \xrightarrow {\operatorname{Homgr} (1 , u)} \operatorname{Homgr} _ {\mathbf {A}} (\mathrm{P}, \mathrm{C}) \xrightarrow {\operatorname{Homgr} (1 , v)} \operatorname{Homgr} _ {\mathbf {A}} (\mathrm{P}, \mathrm{C} ^ {\prime \prime})
\]

和

\[
\operatorname{Homgr} _ {\mathbf {A}} \left(\mathrm{C} ^ {\prime \prime}, \mathrm{E}\right) \xrightarrow {\operatorname{Homgr} (v , 1)} \operatorname{Homgr} _ {\mathbf {A}} (\mathrm{C}, \mathrm{E}) \xrightarrow {\operatorname{Homgr} (u , 1)} \operatorname{Homgr} _ {\mathbf {A}} \left(\mathrm{C} ^ {\prime}, \mathrm{E}\right)
\]

是k-模复形的正合序列。

\begin{enumerate}
\def\labelenumi{\alph{enumi})}
\setcounter{enumi}{1}
\tightlist
\item
  设 \(0 \to \mathbf{C}' \xrightarrow{u} \mathbf{C} \xrightarrow{v} \mathbf{C}'' \to 0\) 是 A-复形的序列，它作为分次 A-模的正合序列是分裂的（例如，当 \(\mathbf{C}'\) 是内射的，或 \(\mathbf{C}''\) 是投射的时，即为这种情况）。则对每个复形 E，序列
\end{enumerate}

\[
0 \rightarrow \operatorname{Homgr} _ {\mathbf {A}} (\mathrm{E}, \mathrm{C} ^ {\prime}) \xrightarrow {\operatorname{Homgr} (1 , u)} \operatorname{Homgr} _ {\mathbf {A}} (\mathrm{E}, \mathrm{C}) \xrightarrow {\operatorname{Homgr} (1 , v)} \operatorname{Homgr} _ {\mathbf {A}} (\mathrm{E}, \mathrm{C} ^ {\prime \prime}) \rightarrow 0
\]

\[
0 \rightarrow \operatorname{Homgr} _ {\mathbf {A}} \left(\mathrm{C} ^ {\prime \prime}, \mathrm{E}\right) \xrightarrow {\operatorname{Homgr} (v , 1)} \operatorname{Homgr} _ {\mathbf {A}} (\mathrm{C}, \mathrm{E}) \xrightarrow {\operatorname{Homgr} (u , 1)} \operatorname{Homgr} _ {\mathbf {A}} \left(\mathrm{C} ^ {\prime}, \mathrm{E}\right)\rightarrow 0
\]

是k-模复形的正合序列。

在情形 \(a\) )，我们注意到序列

\[
\operatorname{Hom} _ {\mathbf {A}} \left(\mathrm{P} _ {p}, \mathrm{C} _ {q} ^ {\prime}\right)\rightarrow \operatorname{Hom} _ {\mathbf {A}} \left(\mathrm{P} _ {p}, \mathrm{C} _ {q}\right)\rightarrow \operatorname{Hom} _ {\mathbf {A}} \left(\mathrm{P} _ {p}, \mathrm{C} _ {q} ^ {\prime \prime}\right)
\]

和

\[
\operatorname{Hom} _ {\mathbf {A}} \left(\mathrm{C} _ {q} ^ {\prime \prime}, \mathrm{E} _ {p}\right)\rightarrow \operatorname{Hom} _ {\mathbf {A}} \left(\mathrm{C} _ {q}, \mathrm{E} _ {p}\right)\rightarrow \operatorname{Hom} _ {\mathbf {A}} \left(\mathrm{C} _ {q} ^ {\prime}, \mathrm{E} _ {p}\right)
\]

对所有 \(p, q \in \mathbf{Z}\) 都是正合的，并且我们应用 II，第 10 页，命题 5 以及 II，第 13 页，命题 7。 \(b\) ) 的证明是类似的。

